\documentclass[11pt]{article}

\usepackage{booktabs}
\usepackage{graphicx}
\usepackage{hyperref}
\usepackage{natbib}
\usepackage[margin=1in]{geometry}

\title{The Paper Crane Benchmark: A Continuously Updated Leaderboard for Protein Binder Design}

\author{%
  Harsha Poonepalle\thanks{Paper Crane. Correspondence: \texttt{harsha@papercrane.bio}.}
  \and
  Vedant Kalipatnapu\footnotemark[1]%
}

\date{\today}

\hypersetup{%
  pdftitle={The Paper Crane Benchmark},
  pdfauthor={Harsha Poonepalle and Vedant Kalipatnapu},
  colorlinks=true,
  linkcolor=blue,
  citecolor=blue,
  urlcolor=blue%
}

\begin{document}

\maketitle

\begin{abstract}
De novo protein binder design has advanced rapidly with the release of tools
such as RFdiffusion~\citep{watson2023rfdiffusion}, ProteinMPNN~\citep{dauparas2022proteinmpnn},
and BindCraft~\citep{pacesa2024bindcraft}, yet the field lacks a shared,
continuously updated benchmark that tracks how these methods perform on a
common set of therapeutic and structural-biology targets under a fixed,
reproducible evaluation harness. We introduce the Paper Crane Benchmark, a
living leaderboard that pins each tool version, target definition, and
structural-validation stack, and recomputes results on a monthly cadence. Each
release is archived as a versioned Zenodo deposition with a citable DOI, a
public dataset of per-design metrics, and a reproducible container image.
Evaluation combines ipSAE~\citep{dunbar2024ipsae}, DockQ~\citep{basu2016dockq},
and self-consistency scores from Boltz-2~\citep{passaro2025boltz2},
AlphaFold~3~\citep{abramson2024alphafold3}, and Chai-1~\citep{chai2024chai1},
with confidence intervals obtained via the bootstrap~\citep{efron1979bootstrap}.
This preprint describes version~\texttt{v\input{data/version.tex}\unskip} of the
benchmark and the infrastructure that keeps it current.
\end{abstract}

\section{Introduction}
\label{sec:introduction}

TODO: motivate the need for a continuously updated protein binder design
benchmark; contrast with static one-shot evaluations; summarize contributions.

\section{Methodology}
\label{sec:methodology}

TODO: describe the evaluation harness, versioning policy, container pinning,
monthly cron cadence, and statistical protocol (bootstrap confidence
intervals~\citep{efron1979bootstrap}).

\section{Targets}
\label{sec:targets}

TODO: enumerate the target set with rationale, PDB IDs, hotspot definitions,
and filter thresholds. Reference per-target metadata shipped in
\texttt{data/targets.json}.

\section{Tools Evaluated}
\label{sec:tools}

TODO: describe each evaluated tool, its pinned version, and its invocation
configuration. Tools currently tracked include
RFdiffusion~\citep{watson2023rfdiffusion},
ProteinMPNN~\citep{dauparas2022proteinmpnn}, and
BindCraft~\citep{pacesa2024bindcraft}, with structural validation by
Boltz-2~\citep{passaro2025boltz2},
AlphaFold~3~\citep{abramson2024alphafold3}, and
Chai-1~\citep{chai2024chai1}.

\section{Results}
\label{sec:results}

\input{tables/results.tex}

TODO: narrate leaderboard results, significance testing, and per-target
breakdowns. Reference \texttt{figures/} for generated plots.

\section{Reproducibility}
\label{sec:reproducibility}

TODO: describe the monthly GitHub Actions workflow, container digests, Zenodo
archival, CITATION.cff versioning, and how external users can reproduce the
reported numbers from a released artifact.

\section{Discussion}
\label{sec:discussion}

TODO: discuss failure modes, calibration of ipSAE~\citep{dunbar2024ipsae} and
DockQ~\citep{basu2016dockq} against wet-lab outcomes, limitations of
structural-only evaluation, and roadmap.

\section{Data Availability}
\label{sec:data}

All raw per-design metrics, target definitions, container digests, and
generated figures and tables for this release are archived on Zenodo with a
citable DOI. The source repository, evaluation harness, and CI workflows are
available at \url{https://github.com/papercrane/mosaic}.

\bibliographystyle{unsrtnat}
\bibliography{references}

\end{document}
